% !TeX root = ../教材题库_全部条目_转写工作稿.tex
% Source: XJ-B1-C03-EX-P075-Y0177
% 第3章 / 3.1 函数 / 教材 p.75 / PDF 第82页
% OCR未校对：公式与文字均可能有误；未完成数学TeX转写；禁止直接发给学生
% 原题及相邻原解截图（核对用）：../题目原图/XJ-B1-C03-EX-P075-Y0177-01.webp
% 本文件只在转写完成后作为题干使用；OCR 条目尚不是可靠题干。
\begin{bankitem}{3.1 例题 例题7}{第3章 / 3.1 函数 / 教材 p.75 / PDF 第82页}{XJ-B1-C03-EX-P075-Y0177}{OCR未校对：公式与文字均可能有误；未完成数学TeX转写；禁止直接发给学生}
\begin{ocrdraft}
[教材p.75 / PDF第82页]
\par
0二7]画出函数不一@的图象，并求不一3），不3）
\par
，不一1)，不1)的值．
\par
解因为丆0）一0
\par
工冫\ensuremath{{}^\circ\mathrm C}，
\par
所以函数丆0）的图象为过原点且平分第一、第二象限的一条
\par
折线，如图3一3．
\par
其中不一3）一3，不3）一3，不一1)一1，丆（1）一1.
\par
由上可知，不一过到也是一个分段函数．
\par
图3．1一3
\par
要解决与分段函数有关的问题，通常要“分段”讨论，
\par
但需注意的是，最后要有一个统一的结论．分段仅仅是用了不同的表达式，分了段
\par
仍然是一个函数．
\end{ocrdraft}
\end{bankitem}
